% !TeX encoding = UTF-8
% !TeX program = pdflatex
% !TeX spellcheck = en_us

\documentclass[uft8,aspectratio=169]{lightweight-slides}

%-----------------------------------------------------------------------

\title{%
    Demo Presentation:\\
    Lightweight Slides Class
}
\author[Author Name]{%
    Author Name\\
    \texttt{author@example.com}%
}
\institute[
    Institution
]{%
    University of Example%
}
\date{%
    \today%
}

\begin{document}
    \nobibliography*

    \frame{\titlepage}

    \begin{frame}
        \frametitle{Table of contents}
        \tableofcontents
    \end{frame}

    \section{Introduction}
    \label{sec:introduction}

    \begin{frame}
        \frametitle{Overview}
        \begin{block}{Blocks and Environments}
            This is a standard block in the lightweight-slides class. It demonstrates the visual hierarchy.
            \begin{itemize}
                \item Bullet point 1
                \item Bullet point 2
            \end{itemize}
        \end{block}
        
        \begin{alertblock}{Alert Block}
            This is an alert block used to emphasize important information.
        \end{alertblock}
    \end{frame}

    \section{Visuals and Features}
    \label{sec:visuals}

    \begin{frame}
        \frametitle{Mathematical Equations}
        The lightweight-slides class supports standard LaTeX math environments:
        \begin{equation*}
            E = mc^2
        \end{equation*}
        \begin{equation*}
            \int_{-\infty}^{\infty} e^{-x^2} dx = \sqrt{\pi}
        \end{equation*}
    \end{frame}

    \begin{frame}
        \frametitle{Tables and Figures}
        \begin{table}[!ht]
            \centering
            \caption{A demo table inside the slides.}
            \begin{tabular}{l c r}
                \hline
                \textbf{Item} & \textbf{Value} & \textbf{Metric} \\
                \hline
                Alpha & 10 & 0.8 \\
                Beta  & 20 & 0.9 \\
                Gamma & 30 & 0.95 \\
                \hline
            \end{tabular}
        \end{table}
    \end{frame}

\end{document}
